\begin{example}[The whole affine scheme from a principal cover]
Suppose $D(f_1),\dots,D(f_r)$ cover $\Spec A$.  Then $(f_1,\dots,f_r)=A$.
The sheaf axiom gives the equalizer
\[
A\longrightarrow\prod_i A_{f_i}
\rightrightarrows
\prod_{i,j}A_{f_if_j}.
\]
Thus the global affine ring is reconstructed from compatible coordinate rings of finitely
many principal affine charts.
\end{example}

\begin{example}[Local coordinates around a prime]
Let $\mfp\in\Spec A$ and let $a/s\in A_{\mfp}$ with $s\notin\mfp$.  Then
$\mfp\in D(s)$ and
\[
D(s)\cong\Spec A_s.
\]
The element $a/s$ is an honest global section on this affine neighborhood.  Every germ on
an affine scheme therefore comes from a section on some affine principal neighborhood.
\end{example}

\begin{example}[The same topology, different affine geometry]
Let $A=k[\varepsilon]/(\varepsilon^2)$ and $B=k$.  Both spectra have one point, so their
underlying spaces are homeomorphic.  But
\[
\Gamma(\Spec A,\OO)=A\not\cong B=\Gamma(\Spec B,\OO),
\]
so Proposition~\ref{prop:vi18-global-recovery} rules out an isomorphism of affine schemes.
\end{example}

\begin{example}[Localization and residue fields]
For $\mfp\in D(f)$, the chart isomorphism gives
\[
(A_f)_{\mfp A_f}\cong A_{\mfp}
\]
and hence the same residue field
\[
\kappa(\mfp A_f)\cong\kappa(\mfp).
\]
Passing to a principal open changes the available global denominators but does not alter
the local geometry at a point that remains in the open.
\end{example}

\begin{example}[Why the projective line obstruction works]
The calculation $\Gamma(\mathbb P^1_k,\OO)=k$ is not merely a scarcity of functions.
For an affine scheme the global ring must reconstruct the entire locally ringed space.
Since $\Spec k$ has one point, a space with projective-line topology cannot be affine with
global section ring $k$.
\end{example}

\begin{example}[Boundary of the chapter]
The ring map $A\to B$ certainly induces a continuous map
\[
\Spec B\longrightarrow\Spec A,\qquad \mfq\longmapsto\varphi^{-1}(\mfq),
\]
but proving that ring maps and affine-scheme morphisms correspond bijectively requires the
sheaf morphism and local-stalk conditions.  That theorem is intentionally VI/19 rather than
being silently absorbed here.
\end{example}
